\chapter{数与数集}\label{ch:g10:numbers}

数学从数开始，但并非所有的数都属于同一类：
计数用的数、负数、分数，以及像 $\sqrt 2$
或 $\pi$ 这样不能用分数表示的数。本章将它们整理为
层层包含的若干类，引入区间来描述数轴上的一段，
并利用绝对值来度量两数之间的距离。

\section{数的分类}

\begin{definition}[数集]\label{def:g10:numbers:sets}
\begin{itemize}
  \item $\N$ 是\emph{自然数}\index{自然数}构成的集合：
        $0, 1, 2, 3, \dots$
  \item $\Z$ 是\emph{整数}\index{整数}构成的集合：$\dots, -2, -1,
        0, 1, 2, \dots$
  \item $\Q$ 是\emph{有理数}\index{有理数}构成的集合：
        所有形如 $\frac{p}{q}$ 的商，其中 $p \in \Z$，$q \in \N$ 且
        $q \neq 0$。
  \item $\R$ 是\emph{实数}\index{实数}构成的集合：所有
        能够放置在数轴上的数。
\end{itemize}
\end{definition}

\begin{notation}
符号 $\in$ 读作“属于”：$3 \in \N$，$-5 \in \Z$，
$\frac23 \in \Q$。符号 $\subset$ 读作“包含于”：每个
自然数都是整数，每个整数都是有理数
（例如 $-5 = \frac{-5}{1}$），每个有理数都是实数，于是
\[
\N \subset \Z \subset \Q \subset \R .
\]
在符号上加斜线表示否定：$\frac12 \notin \Z$。
\end{notation}

\begin{omfigure}
\begin{tikzpicture}[scale=0.9]
  \draw[omExo, thick] (0,0) ellipse (5.4 and 2.9);
  \draw[omProp, thick] (-0.8,0) ellipse (4.0 and 2.2);
  \draw[omThm, thick] (-1.6,0) ellipse (2.6 and 1.5);
  \draw[omDef, thick] (-2.4,0) ellipse (1.3 and 0.8);
  \node[omDef, font=\small\bfseries] at (-2.4,0.45) {$\N$};
  \node[omThm, font=\small\bfseries] at (-1.6,1.1) {$\Z$};
  \node[omProp, font=\small\bfseries] at (-0.8,1.8) {$\Q$};
  \node[omExo, font=\small\bfseries] at (0,2.5) {$\R$};
  \node[font=\small] at (-2.4,-0.15) {$0 \quad 7$};
  \node[font=\small] at (-0.4,-0.9) {$-5$};
  \node[font=\small] at (1.9,-0.5) {$\tfrac23 \quad {-1.7}$};
  \node[font=\small] at (3.9,-0.6) {$\sqrt2 \quad \pi$};
\end{tikzpicture}

{\small 数的层层包含关系：$\N \subset \Z \subset \Q \subset
\R$。每一圈都含有不属于较小一圈的数。}
\end{omfigure}

\begin{example}\label{ex:g10:numbers:classify}
将几个数放入包含它们的最小一类中。
\begin{itemize}
  \item $\frac{28}{4} = 7$，所以 $\frac{28}{4} \in \N$，
        虽然形式上是分数：应当先化简。
  \item $-3.5 = -\frac{35}{10} = -\frac72$ 是有理数，但不是整数。
  \item $0.333\dots$（数字 $3$ 无限重复）等于 $\frac13$，
        是一个有理数。
  \item $\sqrt 2$ 与 $\pi$ 是实数，但不是有理数，其中 $\sqrt 2$ 的情形
        马上就会证明；这类数称为
        \emph{无理数}\index{无理数}。
\end{itemize}
\end{example}

\begin{theorem}[$\sqrt 2$ 是无理数]\label{thm:g10:numbers:sqrt2}
$\sqrt 2$ 不是有理数：不存在平方等于
$2$ 的整数之商。
\end{theorem}

\begin{proof}
采用反证法，分小步进行。
\begin{enumerate}
  \item 假设 $\sqrt 2 = \frac{p}{q}$，其中 $p$ 与 $q$ 为正整数，
        且该分数已经约到最简，因此 $p$ 与 $q$
        不同时为偶数。
  \item 两边平方得 $2 = \frac{p^2}{q^2}$，即
        $p^2 = 2q^2$。因此 $p^2$ 为偶数。
  \item 若 $p$ 为奇数，设 $p = 2k+1$，则
        $p^2 = 4k^2 + 4k + 1$ 为奇数。由于 $p^2$ 为偶数，$p$ 必为
        偶数：$p = 2k$，其中 $k$ 为某个整数。
  \item 代入得 $(2k)^2 = 2q^2$，于是 $4k^2 = 2q^2$，从而
        $q^2 = 2k^2$。仿照第 3 步的论证，$q$ 为偶数。
  \item 此时 $p$ 与 $q$ 都是偶数，与第 1 步矛盾。
        因此假设不成立，$\sqrt 2$ 是无理数。
\end{enumerate}
\end{proof}

\begin{remark}
有理数恰好是小数展开
要么有限结束（如 $\frac72 = 3.5$），要么最终无限重复同一段
（如 $\frac13 = 0.333\dots$ 或 $\frac{1}{7} = 0.142857\,142857\dots$）的数。
像 $\sqrt2 = 1.41421356\dots$ 这样的无理数，其小数数字
永远不会落入循环的模式。这一刻画方法
在本阶段不加证明地承认。
\end{remark}

\section{数轴与区间}

实数充满一条直线：取定原点 $0$ 与单位长度后，
每个实数恰好对应直线上的一个点。

\begin{definition}[区间]\label{def:g10:numbers:interval}
设 $a$、$b$ 为实数，且 $a < b$。
称\emph{区间}\index{区间}为位于两个界之间的
全体实数构成的集合。方括号表示该界包含在内，圆括号表示
该界不包含在内：
\begin{center}
\begin{tabular}{c|c}
记号 & 描述 \\
\hline
$\intcc{a}{b}$ & $a \leq x \leq b$（两端均包含） \\
$\intoo{a}{b}$ & $a < x < b$（两端均不包含） \\
$\intco{a}{b}$ & $a \leq x < b$ \\
$\intoc{a}{b}$ & $a < x \leq b$ \\
$\intco{a}{+\infty}$ & $x \geq a$ \\
$\intoo{-\infty}{b}$ & $x < b$ \\
\end{tabular}
\end{center}
符号 $-\infty$ 与 $+\infty$（读作“无穷”）不是数，只是
一种表示区间无限延伸下去的写法；与它们相邻的括号
一定是圆括号。整条直线 $\R$ 就是区间
$\intoo{-\infty}{+\infty}$。
\end{definition}

\begin{omfigure}
\begin{tikzpicture}[scale=1.0]
  \draw[-Stealth] (-4.4,0) -- (4.6,0) node[below, font=\small] {$x$};
  \foreach \x in {-4,...,4} \draw (\x,-0.07) -- (\x,0.07);
  \foreach \x in {-4,-2,0,2,4}
    \node[below, font=\small] at (\x,-0.1) {$\x$};
  \draw[omDef, very thick] (-3,0.02) -- (1,0.02);
  \fill[omDef] (-3,0.02) circle (2.2pt);
  \draw[omDef, fill=white, thick] (1,0.02) circle (2.2pt);
  \node[omDef, above, font=\small] at (-1,0.15) {$\intco{-3}{1}$};
  \draw[omThm, very thick] (2,0.02) -- (4.5,0.02);
  \draw[omThm, fill=white, thick] (2,0.02) circle (2.2pt);
  \node[omThm, above, font=\small] at (3.4,0.15) {$\intoo{2}{+\infty}$};
\end{tikzpicture}

{\small 数轴上的区间：实心点表示包含的端点，
空心点表示不包含的端点。}
\end{omfigure}

\begin{definition}[交集与并集]\label{def:g10:numbers:interunion}
设 $I$ 和 $J$ 为两个实数集合。
\emph{交集}\index{交集} $I \cap J$（“$I$ 且 $J$”）是
同时属于两者的数构成的集合；
\emph{并集}\index{并集} $I \cup J$（“$I$ 或 $J$”）是
至少属于其中之一的数构成的集合。
\end{definition}

\begin{example}\label{ex:g10:numbers:interunion}
取 $I = \intcc{-1}{3}$ 与 $J = \intoo{1}{5}$。在同一条数轴上画出它们，
两者在 $1$ 与 $3$ 之间有公共部分。由此可得
\[
I \cap J = \intoc{1}{3},
\qquad
I \cup J = \intco{-1}{5}.
\]
注意括号的选取：$1 \notin J$，因此 $1 \notin I \cap J$；而 $3 \in I$ 且
$3 \in J$，因此 $3 \in I \cap J$。
\end{example}

\begin{method}[区间的运算]\label{met:g10:numbers:intervals}
求两个区间的交集或并集，可按以下步骤进行：
\begin{enumerate}
  \item 画出数轴，标出四个端点；
  \item 在数轴上方涂出第一个区间，下方涂出第二个区间；
  \item 两块阴影重叠之处即为交集，并集则是
        至少有一块阴影存在之处；
  \item 判断每个括号时，检验端点本身是否属于
        两个集合（求交集时）或至少一个集合（求并集时）。
\end{enumerate}
\end{method}

\section{绝对值与距离}

\begin{definition}[绝对值]\label{def:g10:numbers:abs}
称实数 $x$ 的\emph{绝对值}\index{绝对值}为
\[
\abs{x} =
\begin{cases}
x & \text{若 } x \geq 0, \\
-x & \text{若 } x < 0.
\end{cases}
\]
例如 $\abs{7} = 7$，$\abs{-4} = 4$。对任意情形，
$\abs{x} \geq 0$。
\end{definition}

\begin{proposition}[数轴上的距离]\label{prop:g10:numbers:distance}
对任意实数 $a$ 与 $b$，数轴上点 $a$ 与
点 $b$ 之间的距离为 $\abs{b - a}$。特别地，$\abs{x}$ 是
$x$ 到 $0$ 的距离。
\end{proposition}

\begin{proof}
若 $b \geq a$，则 $a$ 到 $b$ 的距离为 $b - a \geq 0$，它
等于 $\abs{b-a}$。若 $b < a$，则距离为 $a - b = -(b - a) > 0$，
由绝对值的定义，它仍等于 $\abs{b - a}$。
\end{proof}

\begin{proposition}[绝对值与区间]\label{prop:g10:numbers:absinterval}
设 $a$ 为实数，$r > 0$。则
\[
\abs{x - a} \leq r
\quad\text{当且仅当}\quad
x \in \intcc{a - r}{a + r}.
\]
\end{proposition}

\begin{proof}
$\abs{x-a} \leq r$ 表示 $x$ 到 $a$ 的距离不超过 $r$，
即 $x$ 在任一侧至多偏离 $r$（以 $a$ 为基准），
满足此条件的数恰为位于 $a - r$ 与 $a + r$ 之间的数，
包括两端点。
\end{proof}

\begin{omfigure}
\begin{tikzpicture}[scale=1.1]
  \draw[-Stealth] (-0.8,0) -- (7,0) node[below, font=\small] {$x$};
  \draw[omDef, very thick] (1.5,0.02) -- (5.5,0.02);
  \fill[omDef] (1.5,0.02) circle (2.2pt);
  \fill[omDef] (5.5,0.02) circle (2.2pt);
  \fill (3.5,0) circle (1.6pt);
  \node[below, font=\small] at (3.5,-0.08) {$a$};
  \node[below, font=\small] at (1.5,-0.08) {$a-r$};
  \node[below, font=\small] at (5.5,-0.08) {$a+r$};
  \draw[Stealth-Stealth, omThm] (3.5,0.35) -- (5.5,0.35)
    node[midway, above, font=\small] {$r$};
  \draw[Stealth-Stealth, omThm] (1.5,0.35) -- (3.5,0.35)
    node[midway, above, font=\small] {$r$};
\end{tikzpicture}

{\small 不等式 $\abs{x - a} \leq r$ 描述了距离不超过 $r$ 的数
构成的区间，其中心为 $a$。}
\end{omfigure}

\begin{example}\label{ex:g10:numbers:abs}
解不等式 $\abs{x - 3} \leq 2$。其解为距离不超过 $2$ 的数，
这些数的中心为 $3$，即区间 $\intcc{1}{5}$。反之，区间
$\intcc{-1}{7}$ 的中心为 $\frac{-1+7}{2} = 3$，半径为
$\frac{7-(-1)}{2} = 4$，因此可由 $\abs{x - 3} \leq 4$ 描述。
\end{example}

\section{近似值}

无理数，甚至大多数分数，都不能
精确地写成有限位小数，因此在实际中常用近似值代替它们。

\begin{definition}[给定精度的近似值]\label{def:g10:numbers:approx}
称数 $d$ 为 $x$ 的\emph{近似值}\index{近似值}，精确到
$10^{-n}$ 以内，即 $\abs{x - d} \leq 10^{-n}$。
对小数展开式在第 $n$ 位后截断或四舍五入，均可得到这样的
近似值。
\end{definition}

\begin{example}
由 $\pi = 3.14159\,26\dots$，截断所得的 $3.141$ 与
$3.142$ 都是 $\pi$ 精确到 $10^{-3}$ 以内的近似值。四舍五入
所得值与真值的距离小于 $\frac12 \times 10^{-3}$，而截断只能保证
误差不超过 $10^{-3}$。写出 $3.141 \leq \pi \leq 3.142$，即\emph{夹住}
$\pi$ 位于两个小数界之间。
\end{example}

\section{练习}

\begin{exercise}[$\star$]\label{exo:g10:numbers:1}
对每个数，指出它属于 $\N$、$\Z$、$\Q$、$\R$ 中的
哪一个最小的数集：
\[
\frac{15}{3}, \qquad -7, \qquad \frac{22}{7}, \qquad \sqrt{9}, \qquad
\sqrt{10}, \qquad -2.4, \qquad 0 .
\]
\end{exercise}

\begin{exercise}[$\star$]\label{exo:g10:numbers:2}
用区间记号表示下列各式，再在数轴上
画出：(a) $-2 \leq x < 5$；(b) $x > 3$；(c) $x \leq -1$；
(d) $x$ 与 $2$ 的距离不超过 $3$。
\end{exercise}

\begin{exercise}[$\star$]\label{exo:g10:numbers:3}
计算 $I \cap J$ 与 $I \cup J$，其中
\[
\text{(a) } I = \intcc{-3}{2},\ J = \intco{0}{4};
\qquad
\text{(b) } I = \intoo{-\infty}{1},\ J = \intco{-2}{+\infty}.
\]
\end{exercise}

\begin{exercise}[$\star$]\label{exo:g10:numbers:4}
不用计算器计算：
\[
\abs{-6}, \qquad \abs{4 - 9}, \qquad \abs{-3 - 5}, \qquad
\abs{\sqrt 2 - 1}, \qquad \abs{1 - \sqrt 2}.
\]
\end{exercise}

\begin{exercise}[$\star$]\label{exo:g10:numbers:5}
解下列方程与不等式，并写出解集：
\[
\abs{x} = 5, \qquad \abs{x - 1} = 3, \qquad \abs{x - 4} \leq 1, \qquad
\abs{x + 2} < 3 .
\]
（注意 $\abs{x+2} = \abs{x - (-2)}$ 表示到 $-2$ 的距离。）
\end{exercise}

\begin{exercise}[$\star\star$]\label{exo:g10:numbers:6}
用形如
$\abs{x - a} \leq r$ 或 $\abs{x - a} < r$ 的不等式描述下列各区间：
\[
\intcc{2}{8}, \qquad \intoo{-5}{1}, \qquad \intcc{-7}{-3}.
\]
\end{exercise}

\begin{exercise}[$\star\star$]\label{exo:g10:numbers:7}
证明 $0.272727\dots$（即循环节 $27$ 无限重复）是有理数。
（提示：设 $x$，计算 $100x - x$。）
\end{exercise}

\begin{exercise}[$\star\star$]\label{exo:g10:numbers:8}
判断下列命题的真假，并用论证或反例说明理由。
\begin{enumerate}
  \item 两个整数的和是整数。
  \item 两个整数的商是整数。
  \item 两个有理数的和是有理数。
  \item 一个有理数与一个无理数的和
        是无理数。
\end{enumerate}
\end{exercise}

\begin{exercise}[$\star\star$]\label{exo:g10:numbers:9}
利用 $1.414 \leq \sqrt 2 \leq 1.415$，将数 $2\sqrt2$、
$\sqrt2 + 3$ 与 $-\sqrt2$ 分别用两个小数界夹住。
\end{exercise}

\begin{exercise}[$\star\star\star$]\label{exo:g10:numbers:10}
仿照\cref{thm:g10:numbers:sqrt2}的证明，证明 $\sqrt 3$ 是
无理数。（将“偶数”改为“$3$ 的倍数”，先验证：若
$p^2$ 是 $3$ 的倍数，则 $p$ 也是；
通过余数 $0$、$1$、$2$ 逐一考察（即 $p$ 除以 $3$ 的情形）。）
\end{exercise}

\section{问题：任意两数之间}

\begin{problem}[{周末问题 --- 有理数与无理数交错：
每个区间无论多小，都含有无穷多个有理数与无理数，
而任何测量都永远无法把它们区分开}]
\label{pb:g10:numbers:1}
有理数看上去像一大群（全部分数！），而
无理数则像奇特的例外（$\sqrt2$、$\pi$）。本题
将揭示真实图景：两类数\emph{交错}得如此细密，以至
于数轴上的每个区间，无论
多么微小，都各自含有无穷多个 --- 这带来一些奇特的结论，
例如：无论测量多么精确，
都永远无法判定一个长度是否为有理数。

\medskip
\noindent\textbf{第一部分 --- 四个“王国”。}
\begin{enumerate}
  \item 对每个数，在 $\N$、
        $\Z$、$\Q$、$\R$ 中找出包含它的最小集合
        （见\cref{def:g10:numbers:sets}）：$-7$；\ $\frac{13}{4}$；\
        $\sqrt{16}$；\ $0.121212\ldots$（回忆
        初中分册中的循环小数周末问题）；\ $\sqrt8$；\ $\pi$（承认其
        无理性 ---
        仅在大学分册中证明）。
  \item 证明 $\Q$ 对加法与乘法
        封闭：若 $x = \frac pq$、$y = \frac rs$ 为有理数，
        则将 $x + y$ 与 $xy$ 写成单个分数。
  \item 用反证法推出：（a）有理数与无理数之和
        是无理数；（b）一个\emph{非零}有理数与无理数之积
        是无理数。
  \item 证明无理数集对\emph{两种}运算都不封闭：
        举出两个无理数，其和为有理数；
        再举出两个无理数，其积为有理数。
  \item 指出 $\sqrt2 + \sqrt8$ 与 $\sqrt2 \times \sqrt8$ 分别属于
        四个“王国”中的哪一个（先化简 $\sqrt8$，
        用初中分册的方法）。
\end{enumerate}

\medskip
\noindent\textbf{第二部分 --- 有理数是稠密的。}
\begin{enumerate}[resume]
  \item 求一个严格介于 $3.47$ 与 $3.48$ 之间的有理数；然后
        再求第二个；然后说明如何得到任意多个
        （即初中分册中“没有下一个数”周末问题的“放大”，此时
        已可给出证明）。
  \item 一般定理。设 $a < b$ 为任意两个实数，
        其间距 $g = b - a > 0$。取 $n$ 使得
        $10^{-n} < g$，并考虑 $10^{-n}$ 的倍数
        （步长为 $10^{-n}$ 的小数网格）。说明为什么至少
        有一个格点严格位于 $a$ 与 $b$ 之间，
        并由此得出：\emph{每个正长度的区间
        都含有一个有理数}。
  \item 加强结论：每个这样的区间都含有
        \emph{无穷多个}有理数。（在一个更小的区间内
        再次应用第 7 小题。）
  \item 现在讨论无理数：给定 $a < b$，取有理数 $r$，
        使其严格位于区间内（见第 7 小题），并考虑数
        $r + \frac{\sqrt2}{10^k}$。利用第 3 小题，证明它们
        都是无理数，并且当 $k$ 充分大时它们仍
        位于该区间内：\emph{每个区间也都含有
        无穷多个无理数}。
  \item 解决两个经典谜题：是否存在最小的
        正实数？是否存在一个“恰在
        $3$”之后的实数？用第 7 小题的定理回答这两个问题，并
        向童年时的版本致敬
        （即初中分册中“没有下一个数”周末问题）。
\end{enumerate}

\medskip
\noindent\textbf{第三部分 --- 绝对值，即
$\R$ 的几何。}
\begin{enumerate}[resume]
  \item 解下列各式，并将解表示为区间或区间的并集
        （见\cref{prop:g10:numbers:absinterval}）：
        $\abs{x - 5} = 2$；\quad $\abs{x - 5} < 2$；\quad
        $\abs{x + 1} \geq 3$。
  \item 某活塞的加工尺寸须为 $80$ mm，允许
        误差为 $0.05$ mm。用绝对值写出该要求，
        再用区间表示。两个活塞的测量值分别为
        $79.97$ 与 $80.06$ mm，判断它们是否合格。
  \item 数轴上的三角不等式：
        $\abs{a + b} \leq \abs a + \abs b$。先在
        $(a, b) = (3, -5)$ 与 $(-2, -7)$ 上验证，再证明当 $a$ 与
        $b$ 同号及异号时该式成立，
        并准确说明等号何时成立。
  \item 将方程 $\abs{x - 2} = \abs{x + 4}$ 理解为
        数轴上的距离相等关系，并由此求解。你刚刚求出的是
        初中分册双镜周末问题中的哪一个点，
        降低一维后的对应情形？
  \item 化简 $\sqrt{x^2}$ --- 请务必仔细。检验你的公式
        在 $x = 3$ 与 $x = -3$ 上是否成立，再用它
        借助绝对值求解 $x^2 < 9$。
\end{enumerate}

\medskip
\noindent\textbf{第四部分 --- 任何测量都无法判定的问题。}
\begin{enumerate}[resume]
  \item 一位物理学家测得一根杆子的长度为 $1.41421 \pm 0.00001$ m。
        这样的测量 --- 无论是这一次，还是将来更精确的一次 ---
        能否\emph{证明}这根杆子的长度是无理数，
        或者证明它是有理数？请利用第 7 小题与第 9 小题关于误差区间的结论，
        并说明为什么只有\emph{证明}（如对角线的情形，
        见初中分册中关于无理性的周末问题）
        才能判定无理性。
  \item $1.4$、$1.41$、$1.414$、$1.4142,\dots$ 都是 $\sqrt2$ 的截断值，
        它们都是有理数。它们说明
        $\Q$ 对每个无理数逼近得有多紧？
        请归纳出一般性的结论。
  \item 区间运算：测得的值 $a = 2.5 \pm 0.1$
        与 $b = 1.2 \pm 0.1$。试将 $a + b$ 与 $a \times b$
        夹在某些确定的界之间。哪一种运算会使精度
        降低得更多？
  \item 一条工程规则称，独立测量的误差
        在求和时“相加”。请借助三角不等式（第 13 小题），
        说明为什么 $a + b$ 的误差至多为各误差之和：
        这一不等式正是该条工程规则的
        数学表述。
  \item 结语（短文）：汇总本题及其前身所确立的
        实数轴图景 ——
        没有下一个数
        （初中分册中的周末问题），有理数 $=$
        循环小数（初中分册中循环小数的周末问题），两类数都
        稠密（第 8 小题与第 9 小题），测量永远无法判定
        （第 16 小题）。最后给出本章尚无法证明的一个预告：
        从精确的意义上说，无理数\emph{远远
        多于}有理数 --- 大学分册
        会对无穷进行计数。
\end{enumerate}
\end{problem}
